\section{Results}
\label{sec:results}

All quantitative results in this section come from the frozen
\texttt{controlled\_v3\_real5\_template\_text} run and its provenance-checked
merge. Scores are validation macro-F1. Error bars are sample standard
deviations over seeds 0 and 1 unless a single-run diagnostic is identified.

\subsection{Matched Baselines under Label Skew}

At $\alpha{=}0.1$ and $K{=}5$, local-only training averages 0.297 across the
five clients of one partition. Under the same corpus, partition, encoders,
optimizer, and 24-local-epoch budget, FedAvg reaches
$0.662\!\pm\!0.074$ and FedProx $0.737\!\pm\!0.085$
(Fig.~\ref{fig:baseline}). Their 0.075 gap lies inside the wider sample
standard deviation, so two seeds do not separate them. At $\alpha{=}1$, the
same ordering is 0.523 local-only, $0.820\!\pm\!0.029$ FedAvg, and
$0.862\!\pm\!0.019$ FedProx. Thus, both iterative federated methods sit well
above isolated client training at both skews.

Our SCAFFOLD--AdamW adaptation reaches only $0.070\!\pm\!0.015$ at
$\alpha{=}0.1$ and $0.236\!\pm\!0.205$ at $\alpha{=}1$. Its oscillatory
trajectory points to an unstable AdamW control-variate adaptation; this row
must not be read as a faithful evaluation of classical SGD-based SCAFFOLD.
The FedMME-style one-shot ensemble reaches $0.647\!\pm\!0.080$ at the matched
24-epoch budget and $0.609\!\pm\!0.109$ at its native 100-epoch budget.
Increasing local computation does not repair one-shot aggregation under
severe skew.

\begin{figure}[!t]
\centering
\begin{minipage}[t]{0.465\columnwidth}\centering
\includegraphics[width=\linewidth]{reviewer_matched_results_pa}\\[-5pt]
{\scriptsize\textbf{(a)}}
\end{minipage}\hfill
\begin{minipage}[t]{0.465\columnwidth}\centering
\includegraphics[width=\linewidth]{reviewer_matched_results_pb}\\[-5pt]
{\scriptsize\textbf{(b)}}
\end{minipage}
\caption{Controlled-v3 results at $\alpha{=}0.1$, $K{=}5$:
(a) matched baselines and (b) federated convergence.}
\label{fig:baseline}
\end{figure}

\subsection{Anti-Collapse, Fusion, and Initialization}

The anti-collapse ablation reverses the intended low-$\alpha$ ordering.
With the balanced sampler and entropy-diversity term, F1 is
$0.683\!\pm\!0.089$. Removing the sampler, the diversity term, or both raises
the mean monotonically to $0.716\!\pm\!0.055$,
$0.726\!\pm\!0.035$, and $0.758\!\pm\!0.016$
(Fig.~\ref{fig:anti_missing}(a)). Every one of the eight $\alpha{=}0.1$ runs
finishes at five-class diversity 1.0, including both unregularized seeds, so
final collapse is not observed. The full stack improves the mean minimum
diversity across rounds from 0.70 to 0.90, but it does so with a 0.075 F1
deficit to the unregularized configuration. At $\alpha{=}1$, all four
configurations maintain diversity 1.0 and score 0.809--0.846. We therefore
treat the stack as a transient coverage guard, not an accuracy-improving
component. Likewise, the non-federated early-abort controls score
$0.899\!\pm\!0.022$ with abort disabled and $0.871\!\pm\!0.012$ with it
enabled; neither branch collapses.

\begin{figure}[!t]
\centering
\begin{minipage}[t]{0.435\columnwidth}\centering
\includegraphics[width=\linewidth]{reviewer_matched_results_pc}\\[-5pt]
{\scriptsize\textbf{(a)}}
\end{minipage}\hfill
\begin{minipage}[t]{0.505\columnwidth}\centering
\includegraphics[width=\linewidth]{reviewer_matched_results_pf}\\[-5pt]
{\scriptsize\textbf{(b)}}
\end{minipage}
\caption{Controlled-v3 ablations: (a) anti-collapse F1 and minimum
predicted-class diversity; (b) forced-missing-text performance.}
\label{fig:anti_missing}
\end{figure}

The pooled-data fusion screen is a diagnostic architecture comparison, not an
operational FL result. Across two seeds, its eight operators span 0.750--0.925
F1. The 0.1755 spread is 33 times the 0.0053 median seed standard deviation,
so fusion choice carries real signal in this screen. Residual cross-attention
leads at $0.925\!\pm\!0.001$, while 384-dimensional cross-attention is last at
$0.750\!\pm\!0.002$; with only two seeds, the closely grouped leading methods
remain descriptively rather than statistically ordered. In the separate
operational sweep at $\alpha{=}0.1$, means span 0.636--0.759 but one operator's
sample standard deviation reaches 0.167, so that severe-skew sweep does not
resolve an operational winner (Fig.~\ref{fig:fusion_init}(a)).

The initialization control in Fig.~\ref{fig:fusion_init}(b) closes at
$0.647\!\pm\!0.022$ with all weights random and
$0.831\!\pm\!0.001$ with public encoders and random task heads. A pooled-data
checkpoint opens at 0.874 and closes at $0.877\!\pm\!0.007$, but it was
selected on the reporting split and is a non-deployable diagnostic oracle.
The operational public-encoder initialization retains 94.7\% of its final F1
without accessing pooled training data.

\begin{figure}[!t]
\centering
\begin{minipage}[t]{0.465\columnwidth}\centering
\includegraphics[width=\linewidth]{reviewer_matched_results_pd}\\[-5pt]
{\scriptsize\textbf{(a)}}
\end{minipage}\hfill
\begin{minipage}[t]{0.465\columnwidth}\centering
\includegraphics[width=\linewidth]{reviewer_matched_results_pe}\\[-5pt]
{\scriptsize\textbf{(b)}}
\end{minipage}
\caption{Controlled-v3 comparisons: (a) operational federated fusion under
severe skew; (b) federated initialization and non-federated abort controls.}
\label{fig:fusion_init}
\end{figure}

Under the forced-missing-text metric in Fig.~\ref{fig:anti_missing}(b), zero
filling, deterministic FIN, Gaussian $\beta$-NLL imputation, and
uncertainty-weighted aggregation score $0.688\!\pm\!0.025$,
$0.693\!\pm\!0.004$, $0.706\!\pm\!0.010$, and
$0.706\!\pm\!0.089$. The uncertainty-weighted rule ties the plain
probabilistic rule and is the noisiest, so this matched adaptation does not
show that uncertainty weighting tightens performance.

\subsection{Retrieval-Stage Evaluation}

Exact inner-product FAISS search~\cite{Johnson2021faiss} over training-fitted
TF--IDF note vectors gives condition-macro averages of 0.542 top-1 same-label
accuracy, 0.470 precision@5, and 0.706 mean top-1 cosine similarity across all
600 validation queries (Fig.~\ref{fig:retrieval}). Top-1 accuracy is above the
0.2 label prior but well below the similarity score, consistent with
synthetic templates whose neighbors can be lexically close without sharing a
label. This is a retrieval-stage diagnostic and establishes no clinical
generation quality.

\begin{figure}[!t]
\centering
\includegraphics[width=0.72\columnwidth]{fig_rag_validated}
\caption{Controlled-v3 retrieval metrics by condition over 600 held-out
queries and a 2,400-note index.}
\label{fig:retrieval}
\end{figure}

\subsection{Scalability, Communication, and Resource Cost}

Figure~\ref{fig:scale_accuracy}(a) crosses
$K{=}\{3,5,10,20\}$ with $\alpha{=}\{0.1,1,5\}$. Every
$\alpha{=}5$ cell lies in 0.849--0.916 and every $\alpha{=}1$ cell in
0.818--0.862, whereas $\alpha{=}0.1$ yields 0.643--0.738 with spreads up to
$\pm0.159$. Within a row, increasing $K$ nearly sevenfold moves the mean by
at most 0.10 F1, while changing skew moves it by as much as 0.27. Client count
is therefore secondary to realized label allocation in this grid. Cells show
$A/K$, where $A$ counts shards with at least four samples; the severe
$K{=}20$ split has 17 active clients. Figure~\ref{fig:scale_accuracy}(b)
shows one realized partition in which individual clients hold nearly entire
classes.

\begin{figure}[!t]
\centering
\begin{minipage}[t]{0.465\columnwidth}\centering
\includegraphics[width=\linewidth]{reviewer_systems_scalability_pa}\\[-5pt]
{\scriptsize\textbf{(a)}}
\end{minipage}\hfill
\begin{minipage}[t]{0.465\columnwidth}\centering
\includegraphics[width=\linewidth]{reviewer_systems_scalability_pf}\\[-5pt]
{\scriptsize\textbf{(b)}}
\end{minipage}
\caption{Controlled-v3 scaling: (a) macro-F1 over client count and skew;
(b) one severe-skew client allocation.}
\label{fig:scale_accuracy}
\end{figure}

The GPU-synchronized sections in Fig.~\ref{fig:resources} take 41--62\,s per
round and allocate 4.77--5.92\,GiB on the shared H100. These panels are
execution records, not distributed-throughput curves: clients are simulated
sequentially, and uncontrolled server contention remains.

\begin{figure}[!b]
\centering
\begin{minipage}[t]{0.465\columnwidth}\centering
\includegraphics[width=\linewidth]{reviewer_systems_scalability_pb}\\[-5pt]
{\scriptsize\textbf{(a)}}
\end{minipage}\hfill
\begin{minipage}[t]{0.465\columnwidth}\centering
\includegraphics[width=\linewidth]{reviewer_systems_scalability_pc}\\[-5pt]
{\scriptsize\textbf{(b)}}
\end{minipage}
\caption{Controlled-v3 resource measurements: (a) seconds per round and
(b) peak allocated GPU memory.}
\label{fig:resources}
\end{figure}

For $|\theta|$ trainable parameters, $b$ bytes per parameter, $T$ rounds, and
nominal $K$, the plotted bidirectional volume is
\begin{equation}
V_{\mathrm{nom}}=2KT|\theta|b.
\label{eq:comm}
\end{equation}
Here $|\theta|{=}153{,}935{,}621$, $b{=}4$, and $T{=}8$, giving
27.5/45.9/91.8/183.5\,GiB at $K{=}3/5/10/20$. These are model-state
calculations, not measured network traffic; serialization, compression,
secure aggregation~\cite{Ergun2022globecom}, and algorithm-specific state are
excluded.

The matched single-seed branch audit at $\alpha{=}1$, $K{=}5$ scores text,
image, and concatenated multimodal models at 0.880, 0.768, and 0.830
(Fig.~\ref{fig:costs}(b)). Thus, concatenation improves on image-only by 0.062
but trails text-only by 0.050 in this run. Its one-way FP32 state is 0.573\,GiB
versus 0.248/0.324\,GiB, total timed rounds take 5.6 versus 1.8/4.3 minutes,
and peak allocation is 5.92 versus 2.84/4.33\,GiB. Fusion therefore incurs a
clear systems cost without guaranteeing a gain over the stronger unimodal
branch.

\begin{figure}[!b]
\centering
\begin{minipage}[t]{0.465\columnwidth}\centering
\includegraphics[width=\linewidth]{reviewer_systems_scalability_pd}\\[-5pt]
{\scriptsize\textbf{(a)}}
\end{minipage}\hfill
\begin{minipage}[t]{0.465\columnwidth}\centering
\includegraphics[width=\linewidth]{reviewer_systems_scalability_pe}\\[-5pt]
{\scriptsize\textbf{(b)}}
\end{minipage}
\caption{Controlled-v3 systems cost: (a) formula-derived FP32 model-state
volume; (b) single-seed modality score and measured resource audit.}
\label{fig:costs}
\end{figure}

\section{Discussion and Limitations}

The v3 rerun supports the paper's systems thesis but narrows two model claims.
Federation provides a large advantage over isolated local training, and
partition skew dominates client count. Fusion architecture also matters in
the pooled diagnostic screen. However, the operational concatenated model
does not beat text-only in the matched branch audit, and the anti-collapse
stack does not improve F1 under severe skew. We retain both negative findings
rather than overstate claims contradicted by v3.

The corpus remains a proxy: notes are synthetic and class-paired, not
patient-paired; template wording can leak class information; image class is
partly confounded with source collection; and the split is grouped by neither
patient nor source. Two seeds permit descriptive spreads, not significance,
and deterministic CUDA algorithms were not forced. Sequential single-GPU
simulation omits concurrency, latency, stragglers, dropout, and heterogeneous
hardware. Communication is analytic model-state volume rather than measured
network traffic. FedMME, P-FIN, and SCAFFOLD rows are declared adaptations,
not native replications. Consequently, none of the F1 or retrieval values
estimates clinical safety, generalization, or deployment readiness.

\newpage
\section{Conclusion}

On 3,000 public radiographs with class-paired synthetic notes, matched
controlled-v3 experiments show that FedAvg and FedProx outperform local-only
training at both tested skews, while label skew changes F1 more than a
near-sevenfold change in client count. Public encoders retain 94.7\% of a
pooled diagnostic oracle without pooled operational training. The fusion
screen shows genuine architecture sensitivity, but concatenation trails
text-only in the matched federated branch audit and costs $2.3\times$ its
model state. Removing both anti-collapse components improves severe-skew F1
from 0.683 to 0.758, with all runs finishing at full predicted-class
diversity. These findings define OmniMed-FL as a controlled systems benchmark
and motivate patient-paired, multi-institutional validation with more seeds,
native baseline implementations, and end-to-end network measurements.
